\section{How a perturbation changes the Hamiltonian and states}
\label{app:perturbation}

This appendix makes explicit how the perturbation operator $V$ used in
Sec.~\ref{sec:theory} enters the Hamiltonian and changes its eigenstates,
Majorana wavefunctions, and measured observables.

\paragraph{From a physical parameter to an operator.}
Write the clean Hamiltonian as $H_0$ and the perturbed Hamiltonian as
\begin{equation}
 H'=H_0+V, \qquad V=\sum_\alpha \lambda_\alpha O_\alpha,
 \label{eq:app-general-v}
\end{equation}
where $O_\alpha$ specifies how the environment couples to the chain and
$\lambda_\alpha$ is its strength. For chemical-potential disorder
$\mu_j=\mu_0+\delta\mu_j$, the change relative to Eq.~\eqref{eq:qubit-h} is
\begin{equation}
 V_{\mathrm{dis}}
 =-\frac{1}{2}\sum_j\delta\mu_j Z_j
 =-\sum_j\delta\mu_j\left(n_j-\frac{I}{2}\right),
 \qquad \delta\mu_j\in[-W,W].
 \label{eq:app-vdis}
\end{equation}
The identity contribution has no effect on eigenvectors. If the disorder
coefficient is instead defined as $h_j=-\delta\mu_j$, this becomes
$\sum_j h_j n_j$ up to that constant, which is the shorthand used in the
taxonomy. The boundary poisoning term is instead
\begin{equation}
 V_{\mathrm{poison}}=gX_0=g(c_0+c_0^\dagger).
 \label{eq:app-vpoison}
\end{equation}
Consequently $[V_{\mathrm{dis}},\Pop]=0$, whereas
$\{V_{\mathrm{poison}},\Pop\}=0$.

\paragraph{Change of a many-body wavefunction.}
Let $H_0|n\rangle=E_n|n\rangle$. For an isolated, non-degenerate eigenvalue,
first-order stationary perturbation theory gives
\begin{align}
 E_n'&=E_n+\langle n|V|n\rangle+O(V^2),\\
 |n'\rangle&=|n\rangle+
 \sum_{m\ne n}|m\rangle
 \frac{\langle m|V|n\rangle}{E_n-E_m}+O(V^2).
 \label{eq:app-state-response}
\end{align}
Thus $V$ changes a wavefunction only through matrix elements
$\langle m|V|n\rangle$, while the energy denominators suppress mixing with
states separated by the bulk gap.

\paragraph{Selection by fermion parity.}
For parity eigenstates $\Pop|n,p\rangle=p|n,p\rangle$, with $p=\pm1$, define
$M_{qp}=\langle m,q|V|n,p\rangle$. The commutator or anticommutator with
$\Pop$ implies
\begin{equation}
 \begin{array}{lll}
 [V,\Pop]=0:   &(q-p)M_{qp}=0,&\text{only equal parities can mix},\\
 \{V,\Pop\}=0:&(q+p)M_{qp}=0,&\text{only opposite parities can mix}.
 \end{array}
 \label{eq:app-selection}
\end{equation}
For the nearly degenerate topological states $|e\rangle$ and $|o\rangle$, the
appropriate degenerate-subspace Hamiltonian is
\begin{equation}
 H_{\mathrm{eff}}=
 \begin{pmatrix}
 E_e+v_{ee} & v_{eo}\\
 v_{oe} & E_o+v_{oo}
 \end{pmatrix},
 \qquad v_{pq}=\langle p|V|q\rangle.
 \label{eq:app-heff}
\end{equation}
A parity-even perturbation has $v_{eo}=v_{oe}=0$ and therefore cannot mix the
two sectors. Symmetry alone does not require $v_{ee}=v_{oo}$; locality and the
bulk gap suppress their difference when the two edge modes are well separated.
For a parity-odd perturbation, $v_{ee}=v_{oo}=0$ but the off-diagonal element is
allowed. Its two energies are
\begin{equation}
 E_\pm=\frac{E_e+E_o}{2}
 \pm\sqrt{\left(\frac{E_e-E_o}{2}\right)^2+|v_{eo}|^2}.
 \label{eq:app-splitting}
\end{equation}
Even when the clean finite-size splitting $|E_e-E_o|$ is exponentially small,
a permitted poisoning matrix element can therefore mix and split the sectors.

\paragraph{Majorana profiles and observables.}
For a quadratic, parity-even perturbation, the Bogoliubov--de Gennes (BdG)
matrix changes as $h_{\mathrm{BdG}}'=h_{\mathrm{BdG}}+\delta h$. If
$h_{\mathrm{BdG}}|u_n\rangle=\varepsilon_n|u_n\rangle$ and the level is
isolated, its spatial wavefunction changes by
\begin{equation}
 |\delta u_n\rangle=
 \sum_{m\ne n}|u_m\rangle
 \frac{\langle u_m|\delta h|u_n\rangle}
      {\varepsilon_n-\varepsilon_m}.
 \label{eq:app-bdg-response}
\end{equation}
This is how sub-gap disorder deforms the left- and right-edge Majorana profiles
without necessarily destroying them. The nearly degenerate edge pair itself
must be treated with Eq.~\eqref{eq:app-heff} rather than the non-degenerate
formula.

Finally, for an observable $O$ in a non-degenerate state $|0\rangle$, the
first-order response is
\begin{equation}
 \delta\langle O\rangle=
 2\,\mathrm{Re}\sum_{m\ne0}
 \frac{\langle0|O|m\rangle\langle m|V|0\rangle}{E_0-E_m}+O(V^2).
 \label{eq:app-observable-response}
\end{equation}
Taking $O=\Oedge$ explains why parity-even disorder can continuously deform
the measured edge string even while parity remains exactly fixed. A
parity-odd term instead mixes the two parity sectors through
Eq.~\eqref{eq:app-heff}, so neither $\langle\Pop\rangle$ nor the edge-string
signal is protected.

\section{Fixed-depth MPS memory scaling}
\label{app:mps-memory}

Consider a cut between sites $b$ and $b+1$. Its Schmidt decomposition is
\begin{equation}
 |\psi\rangle=\sum_{\alpha=1}^{\chi_b}s_\alpha
 |\alpha\rangle_{0\ldots b}|\alpha\rangle_{b+1\ldots L-1},
 \label{eq:app-schmidt}
\end{equation}
Here $\chi_b$ is the bond dimension. A product input has $\chi_b=1$, and
single-qubit rotations leave it unchanged. A crossing CNOT has operator
Schmidt rank two,
$\mathrm{CNOT}=|0\rangle\!\langle0|\otimes I
+|1\rangle\!\langle1|\otimes X$, so it at most doubles $\chi_b$.
Since each linear layer crosses the cut once, $r$ repetitions give
$\chi_b\leq2^r$ and hence $\chi=\max_b\chi_b\leq2^r$.

The tensor at site $j$ has dimensions
$\chi_{j-1}\times2\times\chi_j$. The total number of stored complex entries is
therefore bounded by
\begin{equation}
 N_{\mathrm{MPS}}=\sum_{j=0}^{L-1}2\chi_{j-1}\chi_j
 \leq2L\chi^2=O(L\chi^2).
 \label{eq:app-mps-memory}
\end{equation}
For $r=4$, $\chi\leq16$, so the memory is linear in $L$. Kraus trajectories
store pure-state MPS branches rather than the $4^L$ density matrix; general
two-site noise changes the fixed-depth bond constant to $2^{O(r)}$, not the
linear dependence on $L$.
